\documentclass[conference]{IEEEtran}
\usepackage[utf8]{inputenc}
\usepackage[T1]{fontenc}
\usepackage{cite}
\usepackage{amsmath,amssymb}
\usepackage{graphicx}
\usepackage{booktabs}
\usepackage{url}

\title{Design Detectability and a Ceiling Effect in a Preregistered Verifier‑in‑the‑Loop NL\ensuremath{\rightarrow}Config Evaluation}

\author{
\IEEEauthorblockN{Autonomous AI Researcher}
\IEEEauthorblockA{CNSM 2026 Agentic AI Researcher Track}
}

\begin{document}

\maketitle

\begin{abstract}
We report a fully preregistered, archived paired confirmatory evaluation (N = 200 paired multi‑device intents) that measured the marginal benefit of a bounded generate--validate--repair pipeline placing a formal verifier (Batfish) in‑the‑loop versus single‑pass LLM generation for natural‑language\ensuremath{\rightarrow}configuration translation. The preregistered execution used a single shared\_initial\_candidate generation per intent that served as the starting point for both baseline and guarded arms. In the confirmatory cohort the shared candidate satisfied the deterministic validator in every paired episode: baseline\_success\_count = 200, guarded\_success\_count = 200, repair\_attempts = 0/200, paired contingency n11 = 200, n10 = n01 = n00 = 0 (artifact: analysis/paired\_contingency\_table.csv). The exact McNemar two‑sided p = 1.0 and the preregistered 95\% bootstrap percentile CI is degenerate at [0.00, 0.00] (bootstrap\_seed = 1729, resamples = 10000). Because the shared‑candidate pairing prevented the guarded branch from exercising verifier‑guided repairs on failing candidates, the preregistered paired test was uninformative for estimating marginal repair utility. We (1) state the deterministic validator pass rule used, (2) present artifact‑grounded counts and representative diagnostics substantiating the observed ceiling effect, (3) analyze the detectability tradeoff and formal power collapse induced by shared candidate semantics, and (4) give artifact‑grounded design recommendations to restore sensitivity in future preregistered NetOps evaluations.
\end{abstract}

\section{Introduction}
Generative language models are increasingly proposed as components of NetOps toolchains to translate operator intent into device configurations and operational workflows. A prominent mitigation pattern is verifier‑in‑the‑loop: generate a candidate configuration, validate it with a formal network verifier (e.g., Batfish), and, if the verifier reports counterexamples, apply constrained repair iterations. Quantifying the marginal benefit of such guarded pipelines matters for operational adoption because it informs tradeoffs among latency, complexity, and safety. We preregistered study CNSM2026\_GVR\_V1\_001 to answer a confirmatory question: to what extent does an iterative generate--validate--repair loop using a formal verifier reduce critical semantic errors in LLM‑generated multi‑device configurations when evaluated in an emulation‑grounded closed loop? The preregistration fixed a paired design (N = 200 paired intents), stratified sampling, a single shared initial LLM candidate per intent (shared\_initial\_candidate), and at most one verifier‑guided repair iteration.

The run completed as preregistered but produced a design‑induced ceiling effect: the shared initial candidate satisfied the deterministic validator in every paired episode. Consequently, repair logic was never exercised and the preregistered paired contrast (McNemar test) was vacuous. This paper's primary contribution is methodological: using archived artifacts we document the ceiling effect, explain the detectability tradeoff that produced a formal power collapse under the preregistered assumptions, provide representative diagnostics from the archived run, and give concrete, artifact‑grounded recommendations for future preregistered NetOps evaluations that restore sensitivity to marginal verifier utility. We state the study facts plainly early so readers understand the negative/diagnostic character of the result.

\section{Related Work}
Prior work spans intent translation and configuration synthesis, programmatic repair, verifier integration, and agentic orchestration in NetOps. Several recent surveys and architecture proposals argue that verifier‑grounded generate--validate--repair loops can reduce hallucination and semantic errors in NL\ensuremath{\rightarrow}config workflows; executable benchmark efforts and closed‑loop testbeds (for example, NetConfArena and Batfish‑centered evaluations) have demonstrated that formal verifiers detect reachability and policy violations and can improve syntactic/executable correctness in many settings. Concurrent work also documents the execution wall: many benchmarks focus on linguistic or interface tasks while procedural/execution tasks (multi‑device sequence correctness, canary/rollback semantics) remain less standardized. Our study builds on these prior proposals by supplying a fully preregistered, archived paired experiment and by showing how a particular preregistered design choice (shared\_initial\_candidate) can collapse detectability even when the verifier and emulation stack are correctly integrated. The cited record set motivating our design and interpretation includes verifier tool literature and recent NetOps benchmarks and surveys (see cited\_record\_ids).

\section{Methodology}
Preregistration and estimand. The preregistration (CNSM2026\_GVR\_V1\_001) specified the primary estimand paired\_success\_rate\_difference\_guarded\_minus\_baseline, a two‑sided exact McNemar test (alpha = 0.05), and a 95\% bootstrap percentile CI (bootstrap\_seed = 1729, resamples = 10000). Sampling fixed N = 200 confirmatory paired intents stratified into topology clusters: 2 devices (80), 3--4 devices (80), 5--8 devices (40). A reserve of 50 exploratory intents was held aside. Random seed for intent selection: 42. Inclusion/exclusion rules excluded intents referencing proprietary vendor features not reproducible on FRRouting/VyOS.

Generation semantics and orchestration. The run used generation\_semantics = shared\_initial\_candidate: for each intent the pipeline produced a single LLM candidate that served as the starting response for both baseline and guarded conditions. This choice intentionally reduced between‑intent generation variance to isolate repair effects. Declared model: gpt‑5‑mini (execution/model\_configuration.json), provider traces recorded under execution/. Maximum\_model\_calls and repair budgets followed preregistration (maximum\_model\_calls = 400; maximum\_repair\_calls\_per\_task = 1). Orchestration adapter family: hosted\_netops\_gvr\_v1. Emulation used containerized Mininet‑style topologies with FRRouting and VyOS images and Batfish as the formal verifier.

Deterministic validator. We implemented deterministic\_netops\_validator\_v5 as the scorer that maps an LLM candidate to a binary verifier\_pass outcome. The validator assigns verifier\_pass when validation\_after.valid == true and validation\_after.violation\_count == 0. Practically this entailed three sequential checks applied to each candidate: (1) parsing and syntactic executability; (2) task‑payload constraint enforcement (required\_sequence, allowed\_touched\_fields, preserve\_unspecified\_state); and (3) Batfish reachability/policy checks executed against the emulated topology. Validator traces (validation\_before/validation\_after) were archived per candidate in execution/scoring/*.json.

Timeouts, missingness, and contamination controls. Timeouts followed preregistration: LLM call timeout = 60s with one immediate retry; Batfish verification timeout = 120s; emulation timeout = 300s per intent. Exact‑hash duplicate detection and n‑gram overlap flags ran during sampling; exact duplicates were excluded (none found in confirmatory N) and high n‑gram overlap flags were recorded for sensitivity reporting (none flagged). The missingness\_plan deterministically treats technical failures per precommit rules; the run completed without technical failures affecting confirmatory pairs (see execution manifest timestamps and counts below). All raw responses, provider traces, validator traces, and logs were archived for machine verification.

\section{Results}
Execution and accounting. The execution manifest records: started\_at\_utc = 2026‑08‑29T23:01:17.772030+00:00, completed\_at\_utc = 2026‑08‑29T23:16:19.305763+00:00, planned\_episode\_count = 400 (200 baseline + 200 guarded), completed\_episode\_count = 400, failed\_episode\_count = 0, and model\_calls\_used = 200 (execution/execution\_manifest.json; execution/model\_configuration.json). All deterministic consistency checks passed (deterministic\_reconciliation artifact).

Primary artifact‑grounded outcomes. The archived contingency counts were n11 = 200, n10 = 0, n01 = 0, n00 = 0 (analysis/paired\_contingency\_table.csv). Marginal counts: baseline\_success\_count = 200, guarded\_success\_count = 200, giving success rates of 1.0 in both arms (analysis/condition\_summary.csv). The preregistered exact McNemar two‑sided test yields p = 1.0 given zero discordant pairs. The preregistered bootstrap percentile CI (bootstrap\_seed = 1729, resamples = 10000) for the absolute paired difference is degenerate at [0.00, 0.00] (analysis/analysis\_log.jsonl). These numerical outputs are direct reads from the archived analysis artifacts and are reported without modification.

Repair activity and candidate identity. Because of shared\_initial\_candidate semantics, each confirmatory pair used the identical initial LLM candidate for baseline and guarded branches. Representative response files confirm identity between shared and per‑condition outputs (examples: execution/responses/task-000001-shared-initial.txt, execution/responses/task-000001-baseline.txt, execution/responses/task-000001-guarded.txt). For every confirmatory pair validator\_trace shows validation\_after.valid == true, validation\_after.violation\_count == 0, and repair\_applied == false (for example: execution/scoring/task-000001-baseline.json). Consequently, repair\_attempts = 0/200 for the confirmatory cohort: the guarded branch had no opportunity to alter the shared candidate and create discordant outcomes.

Representative diagnostics. To show task complexity despite the ceiling effect, six archived exemplar tasks (execution/representative\_tasks/) summarize difficulty features: changed\_interface\_count values \{1,2,2,2,2,2\} (mean = 1.83), required\_command\_count values \{4,2,6,4,3,6\} (mean = 4.17), dependency\_rule\_count commonly = 2, and difficulty labels spanning medium/hard/very\_hard. These diagnostics demonstrate that confirmatory intents required nontrivial multi‑step sequences and dependency constraints even though the shared candidate satisfied the deterministic validator in every case.

Contamination, exclusions, and missingness. Contamination checks found zero exact‑hash duplicates among confirmatory intents and no high n‑gram overlap flags; no confirmatory episodes were excluded (analysis/contamination\_summary.csv). The missingness summary reports complete\_pairs = 200 and zero failed or unscorable episodes (analysis/missingness\_summary.csv). All such artifacts are available in the archived bundle for independent verification.

Statistical interpretation. The degenerate bootstrap CI and McNemar p = 1.0 do not imply that verifiers are unnecessary in general; they are mathematical consequences of having no discordant pairs when the paired statistic depends solely on discordance. Under shared\_initial\_candidate semantics, the only path to discordance was an applied repair that changed the shared candidate; because every shared candidate passed, repairs were never invoked and the confirmatory contrast could not reveal marginal repair utility.

\section{Discussion}
Detectability tradeoff and formal power collapse. Paired designs reduce between‑intent variance by sharing starting material across conditions. Let S be the event that the shared candidate passes the deterministic validator. When P(S) is close to 1 for the chosen model and generation semantics, the probability of observing any discordant pairs (n10 + n01 > 0) is driven to zero irrespective of conditional repair effectiveness on failing candidates. The preregistered power calculation assumed a nontrivial baseline failure rate and discordant pairs generated by repairs converting failures to passes; those assumptions did not hold in the observed shared candidate distribution, yielding an empirical power collapse. The archived contingency table and validator traces substantiate this formal explanation.

Design recommendations grounded in artifacts. Using only the archived evidence and orchestration semantics, we identify concrete design changes that increase the opportunity for repairs to be exercised without requiring new, post‑lock experiments: (A) independent per‑condition generation: generate separate LLM candidates for baseline and guarded arms (remove shared\_initial\_candidate) so repairs can act on independent failing candidates; (B) multiple independent draws per condition (k > 1) with selection rules informed by validator outcomes to expose diverse failure modes; (C) multi‑model or temperature sampling to increase candidate heterogeneity; (D) inclusion of adversarial/stress strata targeted at ordering and reachability constraints (label these strata exploratory in preregistration); (E) allow extended repair budgets when repair prompts need richer verifier counterexample context. Each change directly addresses the detectability failure observed in the archived artifacts and can be implemented within the same orchestration framework.

Value of trace‑centered evaluation. The run's artifact bundle demonstrates that trace signals (repair\_attempts, repair\_prompt traces, Batfish counterexample provenance, violation\_code taxonomy, repair\_applied flags) are crucial process evidence distinct from binary pass/fail endpoints. Even when final outcomes coincide, these trace metrics reveal whether repair logic was attempted, the nature of counterexamples, and counterexample provenance. Our archived traces (execution/scoring/*.json, execution/provider\_calls/*) enable reproducible secondary analyses focused on process evidence rather than only binary endpoints.

Operational implications. For organizations adopting verifier‑in‑the‑loop deployments, the practical lesson is explicit: if an approval workflow requires a single authoritative initial candidate, a verifier that approves that candidate will not exercise repair logic. To realize repair gains in practice, pipelines should surface diverse candidate proposals, run adversarial perturbations as a canary before approval, or sample multi‑model candidates to increase the likelihood of encountering failing candidates that benefit from repair. The artifacts support these operational recommendations by showing no repair activity under shared\_initial\_candidate semantics despite nontrivial task complexity.

Threats to validity and limits on generalization. Internal validity: the principal internal limitation is the preregistered shared\_initial\_candidate combined with a single model family (gpt‑5‑mini), which restricted within‑study sensitivity. Construct validity: the deterministic validator implements a conservative pass rule (validation\_after.valid == true AND validation\_after.violation\_count == 0) that enforces parsing, task‑constraint, and Batfish checks; alternative validator thresholds would change pass rates. External validity: emulation used FRRouting and VyOS images only---no proprietary vendor images. Conclusion validity: statistical inference collapsed because the required discordant information was absent. These threats and supporting artifacts are documented in the archived bundle (analysis/missingness\_summary.csv; execution/model\_configuration.json).

How future preregistrations can pretest detectability. We recommend preregistrations include a short, automated detectability probe prior to committing the confirmatory sampling plan: draw a small pilot of independent candidates (m = 10--20) under planned generation semantics and compute the pilot pass probability P̂(S). If P̂(S) exceeds a high threshold (e.g., 0.85) for the chosen model and semantics, adapt the design to increase candidate heterogeneity (independent draws, multi‑model sampling, adversarial strata) or declare the confirmatory test exploratory. The present artifacts show P̂(S) = 1.0 under shared\_initial\_candidate for the confirmatory tasks; reporting such a probe before full execution would have signaled the detectability risk.

\section{Limitations}
\begin{itemize}
\item Single model family: only gpt‑5‑mini was evaluated; conclusions do not generalize across models or families.
\item Device and vendor scope: emulation used FRRouting and VyOS images only (no proprietary vendor images).
\item Synthetic deterministic task generator: tasks were deterministically generated and may not capture full operator ambiguity, legacy runbook phrasing, or adversarial intent distributions.
\item Emulator fidelity: containerized Mininet‑style emulation does not fully replicate hardware idiosyncrasies, timing, or race conditions of production devices.
\item Design detectability tradeoff: shared\_initial\_generation was chosen to control generation variance but reduces sensitivity to marginal repair benefits (a principal internal‑validity limitation).
\item Ceiling effect: baseline single‑pass outputs passed the validator for all confirmatory intents, producing a null marginal effect and limiting exploratory inference.
\item Contamination uncertainty: deterministic exact‑hash checks found zero matches in confirmatory N and n‑gram overlap detectors reported no high‑overlap flags, but undetected paraphrase overlap with model training data cannot be fully ruled out.
\item Security scenarios: prompt‑injection, retrieval poisoning, and adversarial telemetry perturbations were not evaluated and remain open threats for agentic NetOps systems.
\end{itemize}

\section{Conclusion}
This fully preregistered, archived paired evaluation exposes a key experimental design lesson for NetOps researchers: generation semantics are a design lever that trades variance control for detectability of verifier utility. In our confirmatory cohort the shared initial candidate matched both baseline and guarded responses for all 200 pairs; no verifier repairs were applied (repair\_attempts = 0/200) and the paired contingency reduced to n11 = 200, rendering McNemar/bootstrap inference vacuous (p = 1.0; 95\% bootstrap CI = [0.00, 0.00]). These statements are artifact‑grounded and reproducible from the bundled execution and analysis artifacts. We recommend that future preregistered NetOps evaluations (1) explicitly assess detectability when selecting generation semantics, (2) include trace‑centered provenance metrics that reveal repair activity, and (3) adopt independent or multiple draws per condition (or multi‑model sampling) when the objective is to quantify marginal repair benefit. The archived artifacts enable community inspection and provide a reproducible record for designing confirmatory NetOps experiments that avoid the ceiling effects documented here.

\section*{Disclosure Statement}
Disclosure Statement (mandatory). This study was executed autonomously from a frozen initial master prompt archived at provenance/master\_prompt.txt (SHA‑256 = 1872df1e\allowbreak{}1805d2d9\allowbreak{}6940456c\allowbreak{}a016bd66\allowbreak{}5d1d5196\allowbreak{}add77f5a\allowbreak{}cdf1582b\allowbreak{}b39b15ba). Declared LLM: gpt‑5‑mini (provider: openai\_responses) as recorded in execution/model\_configuration.json. Orchestration/adapter family: hosted\_netops\_gvr\_v1. Preregistration: CNSM2026\_GVR\_V1\_001 was frozen before execution and specifies the primary estimand, sampling plan (N = 200 confirmatory pairs), shared\_initial\_candidate generation semantics, maximum\_model\_calls = 400, maximum\_repair\_calls\_per\_task = 1, and random seeds (intent selection seed = 42; analysis bootstrap\_seed = 1729). Execution components: containerized Mininet‑style emulation using FRRouting and VyOS images; Batfish served as the formal verifier; deterministic task generator and deterministic\_netops\_validator\_v5 scoring rule were used. Deterministic validator pass rule: verifier\_pass when validation\_after.valid == true AND validation\_after.violation\_count == 0 (validator traces archived under execution/scoring/). The confirmatory run used shared\_initial\_candidate semantics (a single LLM candidate per intent supplied to both baseline and guarded branches) per preregistered design; this choice controlled between‑intent variance but, as reported, made the primary paired test uninformative in this execution. All orchestration, model calls, verifications, emulation runs, scoring, logging, and analysis were performed autonomously per the preregistration; humans acted only as Research Director prior to the autonomy lock and to receive final outputs after autonomous completion. Execution completed without technical restarts affecting confirmatory pairs (completed\_episode\_count = 400; failed\_episode\_count = 0). The artifact bundle includes archived raw responses, validator traces, execution logs, and analysis artifacts (refer to execution/ and analysis/ paths in the bundle) for machine verification.

\begin{thebibliography}{99}
\bibitem{ref1} Matt Brown, Ari Fogel, Daniel Halperin, Victor Heorhiadi, Ratul Mahajan, Todd Millstein, "Lessons from the evolution of the Batfish configuration analysis tool", 2023, 10.1145/3603269.3604866
\bibitem{ref2} Chang Liu, Xiaohui Xie, Xinyi Chen, Yong Cui, "NetConfArena: An Executable Benchmark for LLM Agents in Closed-Loop Network Configuration", 2026, 10.48550/arxiv.2608.23179
\bibitem{ref3} Yuhe Liu, Changhua Pei, Longlong Xu, Bohan Chen, Mingze Sun, Zhirui Zhang, Yongqian Sun, Shenglin Zhang, Kun Wang, Haiming Zhang, Jianhui Li, Gaogang Xie, Xidao Wen, Xiaohui Nie, Ma, Minghua, Pei, Dan, "OpsEval: A Comprehensive IT Operations Benchmark Suite for Large Language Models", 2023, 10.48550/arxiv.2310.07637
\bibitem{ref4} Ghalya Alwhishi, Jamal Bentahar, "A Verifier-Grounded Knowledge-Based LLM Framework with Reinforcement Learning for Reliable NL-to-CTL Specification Generation", 2026, 10.2139/ssrn.6947162
\bibitem{ref5} https://doi.org/10.1145/3603269.3604866
\bibitem{ref6} Jibum Hong, Jibum Hong, Nguyen Van Tu, James Won‐Ki Hong, James Won‐Ki Hong, "A Comprehensive Survey on LLM‐Based Network Management and Operations", 2025, 10.1002/nem.70029
\end{thebibliography}

\end{document}
